% CONTEND - MoSICom 2026 submission draft (IEEE conference format; confirm template + page limit)
% Build on Overleaf (upload this folder) or locally:  latexmk -pdf main.tex
% \TODO{...} marks every number that must come from results/derived/results.json.
\documentclass[conference]{IEEEtran}
\usepackage{cite}
\usepackage{amsmath,amssymb}
\usepackage{graphicx}
\usepackage{booktabs}
\usepackage{xcolor}
\usepackage{url}
\newcommand{\TODO}[1]{\textcolor{red}{[#1]}}
\graphicspath{{../figures/}{figures/}}

\begin{document}

\title{Does Resource Contention Amplify NMS Latency Attacks?\\
A Controlled Factorial Study on a Unified-Memory Edge SoC}

% Remove for double-blind review if the call requires it
\author{\IEEEauthorblockN{\TODO{Author Name}}
\IEEEauthorblockA{\TODO{Affiliation}\\ \TODO{email}}}

\maketitle

\begin{abstract}
Latency attacks inflate the number of candidate boxes entering non-maximum suppression (NMS)
so that an object detector overruns its real-time budget. Prior work has claimed that
interference from co-located workloads magnifies these attacks, but supported the claim only by
comparing different GPUs, and an independent re-evaluation has questioned whether the attacks
matter in practice at all. We test the claim directly on an NVIDIA Jetson Nano, whose CPU and GPU
share one memory system. In a controlled 2$\times$2 factorial design (clean or adversarial input,
with or without a background co-runner), we predict the combined latency from isolated
measurements with a stage-wise proportional-slowdown model and test the prediction with a
pre-registered $\pm$5\% equivalence margin. \TODO{Under memory-bandwidth contention the measured
latency was X\% of the prediction (95\% CI a--b): predictable / amplified.} \TODO{NMS-stage
sensitivity ...} \TODO{At 30 FPS, contention changed the attack's deadline-miss increase from
... to ...} The result bounds the magnification claim to measurable conditions and shows when
isolated profiling is sufficient for real-time budgeting.
\end{abstract}

\begin{IEEEkeywords}
latency attack, non-maximum suppression, resource contention, edge AI, real-time object detection,
availability, Jetson
\end{IEEEkeywords}

\section{Introduction}
Real-time object detectors on edge devices must finish each frame within a fixed budget,
about 33.3\,ms at 30 frames per second. Latency attacks target that budget rather than the
correctness of the prediction: an adversarial input floods NMS with phantom candidate boxes,
and NMS cost grows roughly quadratically with the candidate count
\cite{shapira2023phantom,chen2024overload,wang2021daedalus}.

Two observations motivate this paper. First, the practical size of these attacks is contested:
an independent evaluation across seven platforms and fifteen YOLO models concluded that their
real-world impact is limited under realistic settings \cite{monteuuis2025evade}, and reported
amplification factors depend on configuration details \cite{gu2026survey}. Second, edge
detectors rarely run alone. On integrated system-on-chips the CPU and GPU share DRAM, and
memory-intensive CPU tasks can substantially slow real-time GPU work \cite{ali2018bwlock,jang2020rtod}.

Wang \emph{et al.} \cite{wang2025cantslow} claimed that interference from co-located users
magnifies latency attacks, but their evidence compares a multi-tenant A100 with a single-tenant
RTX~4070Ti, so GPU model and contention change together. A 2026 survey identifies resource
sharing as one of the factors that decide whether a model-level slowdown becomes a system-level
failure, invisible to the model-level measurements most studies report \cite{gu2026survey}.
Whether contention amplifies the attack, or simply slows everything proportionally, has not been
tested under control.

A naive test would not answer the question: if contention slows every stage by a fixed factor
$k$, the additive interaction $T_D-T_B-T_C+T_A=(k-1)(T_B-T_A)$ is positive by construction. We
therefore ask whether adversarial post-processing is \emph{more sensitive} to contention than
clean post-processing, beyond proportional slowdown. Our contributions are:
\begin{itemize}
\item a controlled 2$\times$2 factorial test of the contention-magnification claim on a
unified-memory edge SoC, with two co-runner types that separate memory-bandwidth from CPU contention;
\item a stage-wise proportional-slowdown model that predicts combined latency from isolated
runs, tested with a pre-registered equivalence margin, and stage-level attribution of any deviation;
\item a deadline characterisation at 30\,FPS (absolute miss increase, drops, consecutive misses);
\item an open harness, configuration and raw logs (CONTEND).
\end{itemize}

\section{Background and Related Work}
\textbf{NMS latency attacks.} Daedalus \cite{wang2021daedalus} first broke NMS with dense false
positives; Phantom Sponges \cite{shapira2023phantom} produced a universal perturbation that
overloads NMS while preserving original detections; Overload \cite{chen2024overload} reported
about 10$\times$ inference time on a Jetson NX; Groundswell \cite{xia2026groundswell} improved
stealth and optimisation cost on desktop GPUs. Physical and pipeline-level variants extend the
threat to trackers \cite{ma2024slowtrack,muller2025detstorm,gu2026trackflood,li2026aesop}, and
efficiency attacks on dynamic networks are systematised in \cite{rathnasuriya2025sok,shumailov2021sponge}.

\textbf{Hardware and practicality.} Chen \emph{et al.} attribute the vulnerability to minor page
faults \cite{chen2024resilience}. Wang \emph{et al.} found edge GPUs compute-bound and desktop GPUs
memory-bound under attack and proposed hardware-adaptive adversarial training \cite{wang2025cantslow}.
EVADE \cite{monteuuis2025evade} found the attacks rarely breach 100--500\,ms budgets when NMS runs
on a GPU and that YOLOv8 caps the proposals passed to NMS.

\textbf{Contention on embedded SoCs.} DRAM-bandwidth contention between CPU and GPU on Jetson
boards is well documented for benign workloads \cite{ali2018bwlock,jang2020rtod,valsan2016isolbench}.
To our knowledge, no prior study has combined adversarial NMS workload with controlled contention
on edge hardware (Table~\ref{tab:novelty}).

\begin{table}[t]
\caption{Closest prior work (condensed).}
\label{tab:novelty}
\centering\footnotesize
\begin{tabular}{@{}lccc@{}}
\toprule
Work & Controlled contention & Joint 2$\times$2 & Edge HW \\
\midrule
\cite{wang2025cantslow} & no (cross-GPU) & no & partly \\
\cite{monteuuis2025evade} & no & no & yes \\
\cite{chen2024overload,chen2024resilience} & no & no & yes \\
\cite{ali2018bwlock} & yes & clean only & yes \\
This work & yes & yes & yes \\
\bottomrule
\end{tabular}
\end{table}

\section{Threat and System Model}
The attacker controls camera input pixels through a universal perturbation (digital, as in
\cite{shapira2023phantom}) and does not control the platform. The detector runs on an
NVIDIA Jetson Nano (quad-core Cortex-A57, 128-core Maxwell GPU, 4\,GB LPDDR4 shared by CPU and GPU)
alongside other on-board software, modelled by a background co-runner pinned to the remaining
CPU cores. Frames arrive at \TODO{30} FPS; the per-frame deadline is \TODO{33.3}\,ms.

\section{Method}
\subsection{Factorial design}
Cells A (clean, no load), B (adversarial, no load), C (clean, load) and D (adversarial, load)
are run in a seeded random order, \TODO{n} independent runs per cell, each \TODO{3000} frames after
300 warm-up frames, with cool-downs between runs. The run, not the frame, is the experimental unit.
Two contention types are tested: a memory-bandwidth co-runner streaming over a 256\,MB buffer,
and a cache-resident integer co-runner that loads the CPU without DRAM traffic (control).

\subsection{Stage-wise proportional-slowdown model}
End-to-end service time is the sum of stage times $T=\sum_s T_s$ for
$s\in\{\text{pre-processing, inference, NMS, copy}\}$. Contention sensitivity is measured on
clean input and the combined latency is predicted from isolated cells only:
\begin{equation}
k_s=\frac{\bar T_s(C)}{\bar T_s(A)},\qquad \hat T_D=\sum_s k_s\,\bar T_s(B).
\end{equation}
The test statistic is $R=\log(\bar T_D/\hat T_D)$, with a run-level bootstrap \cite{efron1993bootstrap}.
H0 (predictable) holds when the $(1-2\alpha')$ CI of $R$ lies inside $\pm\log(1.05)$ (two one-sided
tests \cite{lakens2017tost}); H1 (amplified) when the $(1-\alpha')$ CI lies above $+\log(1.05)$;
$\alpha'=0.05$ divided by the number of contention types. For RQ2 we compare
$k_s^{\text{adv}}=\bar T_s(D)/\bar T_s(B)$ with $k_s^{\text{clean}}$ per stage.

\subsection{Deadline metrics}
In open-loop runs frames are released at a fixed rate into a bounded drop-oldest queue of length 4.
We report the deadline-miss rate (drops count as misses), the absolute miss increase
$P_{\text{miss}}(D)-P_{\text{miss}}(C)$ versus $P_{\text{miss}}(B)-P_{\text{miss}}(A)$, maximum
consecutive misses and P95 tardiness, following \cite{gu2026survey}.

\section{Experimental Setup}
JetPack \TODO{4.6.x} (TensorRT \TODO{8.x}, CUDA 10.2), MAXN power mode, locked clocks, fan at maximum.
Model \TODO{YOLOv8n} at \TODO{416}$\times$\TODO{416}, FP16 TensorRT engine, chosen so that clean
latency is 50--70\% of the deadline. NMS mirrors Ultralytics defaults \cite{jocher2023ultralytics}
(confidence 0.25, IoU 0.7, max\_det 300) and runs on the GPU. Perturbation: Phantom Sponges
objective \cite{shapira2023phantom}, $L_\infty\le16/255$, optimised on 200 COCO val2017 images
\cite{lin2014coco} disjoint from the 500-image evaluation subset. Gate checks: candidates increased
\TODO{x}$\times$; each co-runner slowed clean latency by \TODO{y}\%.

\section{Results}
\subsection{RQ1: Is the combined effect predictable?}
\begin{figure}[t]\centering
\includegraphics[width=\columnwidth]{fig2_rq1_prediction.pdf}
\caption{Measured over predicted latency of cell D per contention type, 95\% CI; grey band: $\pm$5\% equivalence margin.}
\label{fig:rq1}\end{figure}
\TODO{Fill from results.json: ratio, CI, decision per contention type; report the naive interaction $I_T$ for contrast.}

\subsection{RQ2: Which stage carries the effect?}
\begin{figure}[t]\centering
\includegraphics[width=\columnwidth]{fig3_rq2_stage_sensitivity.pdf}
\caption{Per-stage contention slowdown for clean and adversarial input (95\% CI).}
\label{fig:rq2}\end{figure}
\TODO{NMS-stage $k^{\text{adv}}/k^{\text{clean}}$ with CI; memory-bandwidth vs CPU co-runner.}

\subsection{RQ3: Deadline consequences}
\begin{figure}[t]\centering
\includegraphics[width=\columnwidth]{fig4_rq3_trace.pdf}
\caption{Per-frame latency in open-loop runs under contention, clean vs adversarial.}
\label{fig:rq3}\end{figure}
\TODO{Table of miss rate, drops, max consecutive misses per cell.}

\section{Discussion and Limitations}
\TODO{What the result means for \cite{wang2025cantslow} and \cite{monteuuis2025evade}.}
Threats to validity: one platform and one detector family; a digital perturbation; synthetic
co-runners; \TODO{n} runs per cell under a compressed schedule. Results are bounded to the tested
configuration.

\section{Conclusion}
\TODO{One paragraph.} Future work: additional platforms (including CPU-only boards), tracker stages
and contention-aware admission caps.

\bibliographystyle{IEEEtran}
\bibliography{refs}
\end{document}
